\documentclass[a4paper,11pt]{article}
\usepackage[T2A]{fontenc}
\usepackage[utf8]{inputenc}
\usepackage[russian]{babel}
\usepackage{amsmath,amssymb}
\usepackage{geometry}
\usepackage{booktabs,tabularx,array}
\usepackage{microtype}
\usepackage{helvet}
\renewcommand{\familydefault}{\sfdefault}
\usepackage{xcolor}
\usepackage{tikz}
\usepackage{hyperref}
\usepackage{fancyhdr}
\usepackage{enumitem}
\usepackage{icomma}
\geometry{left=23mm,right=22mm,top=22mm,bottom=25mm}
\definecolor{accent}{HTML}{136D70}
\definecolor{accentdark}{HTML}{164E51}
\definecolor{quiet}{HTML}{526064}
\hypersetup{colorlinks=true,linkcolor=accentdark,pdfauthor={Лёвин Артём Александрович},pdftitle={Ревью домашней работы — Ксения Васильченко}}
\setlength{\parindent}{0pt}
\setlength{\parskip}{0.52em}
\setlength{\headheight}{15pt}
\setlength{\tabcolsep}{5pt}
\setlist[itemize]{leftmargin=1.5em,itemsep=.2em,topsep=.2em}
\pagestyle{fancy}
\fancyhf{}
\fancyhead[L]{\small\textcolor{quiet}{Математика · Ревью домашней работы}}
\fancyhead[R]{\small\textcolor{quiet}{04.11.2026}}
\fancyfoot[C]{\footnotesize\textcolor{quiet}{Лёвин Артём Александрович эксклюзивно для Ксении Васильченко}}
\fancyfoot[R]{\footnotesize\thepage}
\renewcommand{\headrulewidth}{0pt}
\renewcommand{\footrulewidth}{0pt}
\newcommand{\taskhead}[2]{\phantomsection\label{task:#1}{\LARGE\color{accentdark}\bfseries Задание №\,#1\par}\vspace{1mm}{\small\color{quiet}#2\par}\vspace{3mm}}
\newcommand{\parttitle}[1]{\par\vspace{1.6mm}\noindent{\normalsize\color{accentdark}\bfseries #1}\par\vspace{-1.5mm}}
\begin{document}
\thispagestyle{empty}
\vspace*{39mm}
\begin{center}
{\small\color{accent}МАТЕМАТИКА}\par
\vspace{12mm}
{\Huge\bfseries Ревью\\ домашней работы}\par
\vspace{9mm}
{\large Дроби, части числа и периметр прямоугольника}\par
\vspace{19mm}
{\Large Ксения Васильченко}\par
\vspace{18mm}
{\normalsize Лёвин Артём Александрович}\par
\vspace{3mm}
{\normalsize 4 ноября 2026 года}\par
\end{center}
\vfill
\begin{center}
{\small\color{quiet}Лёвин Артём Александрович эксклюзивно для Ксении Васильченко}\par
\vspace{8mm}
\begin{tikzpicture}
\draw[accent!55,line width=0.75pt] (-5,0) .. controls (-3,0.55) and (-1,-0.55) .. (0,0)
 .. controls (1,0.55) and (3,-0.55) .. (5,0);
\end{tikzpicture}
\end{center}
\newpage
\section*{Сводная таблица}
В таблицу включены только задания, требующие исправления или завершения. Нажмите на номер, чтобы перейти к подробному разбору.
\par\vspace{5mm}
\noindent{\small
\renewcommand{\arraystretch}{1.55}
\begin{tabularx}{\linewidth}{@{}p{10mm}p{25mm}X p{21mm}p{20mm}@{}}
\toprule
№ & Тема & Суть ошибки & Ваш ответ & \shortstack[l]{Верный\\ответ} \\
\midrule
\hyperref[task:3]{3} & Число по его части & Перепутаны делимое и делитель. & $\frac{1}{35}$ & $35$ \\
\hyperref[task:4]{4} & Число по его части & Перепутан порядок деления. & $\frac{2}{90}$ & $45$ \\
\hyperref[task:10]{10} & Прямоугольник & Формулы записаны для квадрата; периметр не найден. & Не указан & $18\frac12$ см \\
\bottomrule
\end{tabularx}
}
\vfill
{\color{quiet}\small Главная тема для повторения: если известна часть числа, всё число находят делением известной части на долю, выраженную дробью или десятичным числом.}
\newpage
\thispagestyle{empty}
\null\vfill
\begin{center}
{\Huge\bfseries\color{accentdark}Разбор ошибок}
\end{center}
\vfill\null
\newpage
\taskhead{3}{Нахождение числа по известной дробной части}
\parttitle{Текст задания}
Найдите число, если $\frac37$ этого числа равны $15$.
\parttitle{Ваше решение и ответ}
\[
\frac37:15=\frac37:\frac{15}{1}=\frac37\cdot\frac1{15}=\frac1{35}.
\]
Ваш ответ: $\frac1{35}$.
\parttitle{Ошибка и её точное место}
\textcolor{red}{\textbf{Ошибка:}} при переходе от условия к действию на деление перепутан порядок чисел. Вы вычислили $\frac37:15$, хотя нужно вычислить $15:\frac37$. Выбранное Вами деление выполнено верно; именно выбор действия для данной задачи неверен.
\parttitle{Последний правильный шаг и первый неверный}
Из условия следует, что неизвестное число $x$ удовлетворяет равенству
\[\frac37\cdot x=15.\]
Это исходное правильное равенство. Первый неверный шаг в Вашей записи --- $\frac37:15$ вместо $15:\frac37$.
\parttitle{Почему это неверно}
Фраза «$\frac37$ числа равны 15» означает, что число сначала \emph{умножили} на $\frac37$ и получили $15$. Чтобы восстановить исходное число, нужно отменить умножение делением на $\frac37$. При делении порядок важен: $a:b$ обычно не равно $b:a$. Кроме того, искомое число должно быть больше $15$, ведь $\frac37$ --- часть целого, меньшая единицы.
\parttitle{Ключевая идея}
\textcolor{blue!60!black}{\textbf{Ключевая идея:}} если $\frac{a}{b}$ числа равны $c$, то всё число равно $c:\frac{a}{b}$.
\parttitle{Исправление и полное решение}
\begin{align*}
\frac37x&=15,\\
x&=15:\frac37,\\
x&=15\cdot\frac73,\\
x&=\frac{15\cdot7}{3}=5\cdot7=35.
\end{align*}
\parttitle{Проверка}
$\frac37\cdot35=\frac{105}{7}=15$. Получено именно число из условия.
\parttitle{Итог}
\textcolor{teal!60!black}{\textbf{Ответ:}} $35$.
\parttitle{Задача на закрепление}
Найдите число, если $\frac58$ этого числа равны $30$.
\newpage
\taskhead{4}{Нахождение числа по десятичной части}
\parttitle{Текст задания}
Найдите число, если $0{,}4$ этого числа равны $18$.
\parttitle{Ваше решение и ответ}
В тетради выполнено деление
\[
0{,}4:18=\frac4{10}:\frac{18}1
=\frac4{10}\cdot\frac1{18}=\frac2{90}.
\]
Ваш ответ: $\frac2{90}$.
\parttitle{Ошибка и её точное место}
\textcolor{red}{\textbf{Ошибка:}} при составлении первого действия записано $0{,}4:18$. Это действие находит частное чисел $0{,}4$ и $18$, но не восстанавливает всё число по его части.
\parttitle{Последний правильный шаг и первый неверный}
Перевод $0{,}4=\frac4{10}$ выполнен верно; деление $0{,}4:18$ само по себе тоже рассчитано правильно. Первое неверное решение для условия задачи --- выбор порядка деления: вместо $18:0{,}4$ записано $0{,}4:18$.
\parttitle{Почему это неверно}
Если $0{,}4$ неизвестного числа равны $18$, значит неизвестное число умножили на $0{,}4$ и получили $18$. Число $0{,}4$ меньше единицы, поэтому исходное число должно быть больше $18$. Ответ $\frac2{90}$ гораздо меньше $18$; он противоречит смыслу условия.
\parttitle{Ключевая идея}
\textcolor{blue!60!black}{\textbf{Ключевая идея:}} число по его части находят делением известной части на соответствующую долю.
\parttitle{Исправление и полное решение}
Обозначим всё число через $x$. Тогда
\begin{align*}
0{,}4x&=18,\\
x&=18:0{,}4,\\
x&=18:\frac4{10},\\
x&=18\cdot\frac{10}{4},\\
x&=\frac{180}{4}=45.
\end{align*}
\parttitle{Проверка}
$45\cdot0{,}4=18$. Условие выполнено.
\parttitle{Итог}
\textcolor{teal!60!black}{\textbf{Ответ:}} $45$.
\parttitle{Задача на закрепление}
Найдите число, если $0{,}6$ этого числа равны $27$.
\newpage
\taskhead{10}{Площадь и периметр прямоугольника}
\begingroup\setlength{\parskip}{0.25em}\small
\parttitle{Текст задания}
Площадь прямоугольника равна $15\frac34$ см$^2$, а одна сторона --- $2\frac14$ см. Найдите периметр прямоугольника.
\parttitle{Ваше решение и ответ}
В тетради записаны формулы $P=(a+a)\cdot2$ и $S=a\cdot a$. Далее Вы правильно разделили площадь на известную сторону:
\[
15\frac34:2\frac14=\frac{63}{4}:\frac94
=\frac{63}{4}\cdot\frac49=7.
\]
После этого стоит «2) ?». Итоговый периметр не вычислен, ответ отсутствует.
\parttitle{Ошибка и её точное место}
\textcolor{red}{\textbf{Ошибка:}} формулы $P=2(a+a)$ и $S=a\cdot a$ применимы к квадрату с равными сторонами. В условии дан прямоугольник; его стороны могут быть различны. Кроме того, после правильного нахождения второй стороны решение осталось незавершённым.
\parttitle{Последний правильный шаг и первый неверный}
Выполненное деление $\frac{63}{4}:\frac94=7$ --- правильный шаг, его можно сохранить. Ошибка в формулах возникла раньше: для прямоугольника необходимо использовать две стороны $a$ и $b$; после получения $7$ см нужно было найти сумму длин всех четырёх сторон.
\parttitle{Почему это неверно}
Площадь прямоугольника равна произведению \emph{длины и ширины}: $S=a\cdot b$. Периметр --- сумма всех сторон: $P=a+b+a+b=2(a+b)$. Формулы с $a\cdot a$ и $2(a+a)$ подходят лишь для частного случая $a=b$. Здесь найденная вторая сторона $7$ см отличается от известной стороны $2\frac14$ см.
\parttitle{Ключевая идея}
\textcolor{blue!60!black}{\textbf{Ключевая идея:}} сначала найдите неизвестную сторону по формуле $b=S:a$, затем найдите периметр по формуле $P=2(a+b)$.
\parttitle{Исправление и полное решение}
Пусть $a=2\frac14$ см, $S=15\frac34$ см$^2$. Переведём смешанные числа в неправильные дроби:
\[
a=\frac94\text{ см},\qquad S=\frac{63}{4}\text{ см}^2.\]
Найдём вторую сторону:
\[
b=S:a=\frac{63}{4}:\frac94
=\frac{63}{4}\cdot\frac49=7\text{ см}.\]
Теперь вычислим периметр:
\begin{align*}
P&=2(a+b),\\
P&=2\left(2\frac14+7\right)
=2\left(9\frac14\right)
=18\frac12\text{ см}.
\end{align*}
\parttitle{Проверка}
Площадь: $7\cdot2\frac14=7\cdot\frac94=\frac{63}{4}=15\frac34$ см$^2$. Периметр: $7+2\frac14+7+2\frac14=18\frac12$ см.
\parttitle{Итог}
\textcolor{teal!60!black}{\textbf{Ответ:}} $18\frac12$ см.
\parttitle{Задача на закрепление}
Площадь прямоугольника равна $13\frac12$ см$^2$, а одна его сторона --- $2\frac14$ см. Найдите периметр прямоугольника.
\endgroup
\end{document}
